\documentclass[12pt,a4paper]{article}
\usepackage[utf8]{inputenc}
\usepackage[spanish]{babel}
\usepackage{amsmath}
\usepackage{amsfonts}
\usepackage{amssymb}
\usepackage[left=2cm,right=2cm,top=2cm,bottom=2cm]{geometry}

\begin{document}

\begin{center}
    \Large \textbf{Universidad Nacional de Ingeniería (UNI)} \\
    \large \textbf{Programa de Ingeniería de Sistemas} \\
    \vspace{0.5cm}
    \textbf{Prueba Diagnóstica: Pre-cálculo e Introducción a Límites} \\
    \textbf{Asignatura: Matemática I}
\end{center}

\vspace{0.5cm}
\noindent\textbf{Nombre:} \rule{8cm}{0.4pt} \hfill \textbf{Fecha:} \rule{3cm}{0.4pt}
\vspace{0.5cm}

\section*{Parte I: Fundamentos Algebraicos}
\textit{(Estos ejercicios evalúan las herramientas necesarias para manipular formas indeterminadas).}

\begin{enumerate}
    \item \textbf{Simplificación:} Simplifica la siguiente expresión algebraica, indicando si existe alguna restricción para el valor de $x$:
    \[ \frac{x^2 - 4x + 3}{x - 3} \]

    \vspace{2.5cm}

    \item \textbf{Racionalización:} Racionaliza el numerador y simplifica la siguiente expresión:
    \[ \frac{\sqrt{x+4} - 2}{x} \]
\end{enumerate}

\vspace{2.5cm}

\section*{Parte II: Comprensión de Funciones}
\textit{(Para entender el comportamiento de las funciones antes de aplicarles el operador límite).}

\begin{enumerate}
    \setcounter{enumi}{2}
    \item \textbf{Dominio:} Determina el dominio de la función $f(x) = \frac{5}{\sqrt{x^2 - 16}}$.

    \vspace{2.5cm}

    \item \textbf{Evaluación y Tasa de Cambio:} Dada la función $f(x) = 2x^2 - x$, encuentra y simplifica la siguiente expresión:
    \[ \frac{f(x+h) - f(x)}{h} \]
\end{enumerate}

\vspace{2.5cm}

\section*{Parte III: Noción Intuitiva de Límite}
\textit{(Exploración numérica y conceptual).}

\begin{enumerate}
    \setcounter{enumi}{4}
    \item \textbf{Aproximación Numérica:} Considera la función $g(x) = \frac{\sin(x)}{x}$ (asumiendo que $x$ está en radianes). 
    \begin{itemize}
        \item Usando una calculadora, evalúa la función para $x = 0.1$, $x = 0.01$ y $x = 0.001$. 
        \item ¿A qué valor parece acercarse $g(x)$ cuando $x$ se acerca a $0$? 
        \item ¿Qué ocurre matemáticamente si intentas evaluar $g(0)$ directamente?
    \end{itemize}

    \vspace{3cm}

    \item \textbf{Aplicación Práctica:} El tiempo de respuesta $T$ (en milisegundos) de un servidor web al recibir $n$ miles de peticiones simultáneas está modelado por la función:
    \[ T(n) = \frac{50n}{n+2} \]
    ¿Qué sucede con el tiempo de respuesta del servidor a medida que la cantidad de peticiones crece indefinidamente (es decir, cuando $n$ toma valores cada vez más grandes)?
\end{enumerate}

\end{document}